\section{Introduction}

The United States interstate highway system is the largest infrastructure investment in American history. At 48,800 miles, it connects every major metropolitan area and carries 72\% of the nation's freight value by truck \citep{BTS2024}. Yet 20.4\% of Americans --- 66.5 million people --- live in counties whose population centroid is more than 30 miles from any interstate on-ramp.

This is not a uniform distribution. The gaps cluster in specific geographic zones: the entire northern tier of the country from the Canadian border down to I-90 and I-94, a 300--500 mile band where no interstate operates; the central Appalachian region between I-40 and I-70 where West Virginia, eastern Kentucky, and western Virginia have no east-west interstate crossing; the rural Gulf South between I-10 and I-20 across Louisiana, Mississippi, and Arkansas; and the Rural West, where enormous Nevada, Utah, and New Mexico counties have populations hundreds of miles from any on-ramp.

The 30-mile standard is not arbitrary. It corresponds to approximately 35 minutes of rural driving --- the threshold below which emergency medical response times become survivable, the range within which agricultural freight moves cost-effectively during harvest windows, and the distance at which interstate access supports rather than constrains rural economic opportunity \citep{USDA_ERS2020}. Communities beyond this threshold are systematically disadvantaged in healthcare outcomes, agricultural productivity, labor market access, and disaster evacuation capacity.

This paper makes four contributions. First, we present the first systematic population-weighted coverage analysis of the US interstate network at the county level, using Census Bureau Gazetteer county centroids and TIGER/Line highway geometry. Second, we classify the 1,510 continental US gap counties into four types based on geographic clustering, corridor availability, and population density. Third, we rank 12 proposed interstate corridors by coverage efficiency: the ratio of gap-county population served to estimated construction cost. Fourth, we demonstrate that gap counties score 41\% higher than non-gap counties on the C3 (Economic Opportunity Access) dimension of the ROUTE scoring framework --- confirming that the coverage gap and the economic opportunity gap are the same gap.

The findings directly inform Interstate 2.0 investment sequencing. Closing the coverage gap is not a question of adding lanes to existing congested corridors; it is a question of building into geographic zones the system never reached.

\subsection*{Contributions}

\begin{enumerate}
\item First population-weighted county-level coverage analysis of the US interstate network: 66.5 million Americans (20.4\%) beyond the 30-mile access threshold
\item Four-type gap taxonomy with geographic cluster identification and quantified population exposure
\item Ranking of 12 proposed corridors by coverage efficiency (gap population per \$B of investment)
\item Empirical confirmation that coverage gaps and economic opportunity gaps are geographically co-located
\end{enumerate}
